% !TEX root = paper.tex

%% ODER: format ==         = "\mathrel{==}"
%% ODER: format /=         = "\neq "
%
%
\makeatletter
\@ifundefined{lhs2tex.lhs2tex.sty.read}%
  {\@namedef{lhs2tex.lhs2tex.sty.read}{}%
   \newcommand\SkipToFmtEnd{}%
   \newcommand\EndFmtInput{}%
   \long\def\SkipToFmtEnd#1\EndFmtInput{}%
  }\SkipToFmtEnd

\newcommand\ReadOnlyOnce[1]{\@ifundefined{#1}{\@namedef{#1}{}}\SkipToFmtEnd}
\usepackage{amstext}
\usepackage{amssymb}
\usepackage{stmaryrd}
\DeclareFontFamily{OT1}{cmtex}{}
\DeclareFontShape{OT1}{cmtex}{m}{n}
  {<5><6><7><8>cmtex8
   <9>cmtex9
   <10><10.95><12><14.4><17.28><20.74><24.88>cmtex10}{}
\DeclareFontShape{OT1}{cmtex}{m}{it}
  {<-> ssub * cmtt/m/it}{}
\newcommand{\texfamily}{\fontfamily{cmtex}\selectfont}
\DeclareFontShape{OT1}{cmtt}{bx}{n}
  {<5><6><7><8>cmtt8
   <9>cmbtt9
   <10><10.95><12><14.4><17.28><20.74><24.88>cmbtt10}{}
\DeclareFontShape{OT1}{cmtex}{bx}{n}
  {<-> ssub * cmtt/bx/n}{}
\newcommand{\tex}[1]{\text{\texfamily#1}}	% NEU

\newcommand{\Sp}{\hskip.33334em\relax}


\newcommand{\Conid}[1]{\mathit{#1}}
\newcommand{\Varid}[1]{\mathit{#1}}
\newcommand{\anonymous}{\kern0.06em \vbox{\hrule\@width.5em}}
\newcommand{\plus}{\mathbin{+\!\!\!+}}
\newcommand{\bind}{\mathbin{>\!\!\!>\mkern-6.7mu=}}
\newcommand{\rbind}{\mathbin{=\mkern-6.7mu<\!\!\!<}}% suggested by Neil Mitchell
\newcommand{\sequ}{\mathbin{>\!\!\!>}}
\renewcommand{\leq}{\leqslant}
\renewcommand{\geq}{\geqslant}
\usepackage{polytable}

%mathindent has to be defined
\@ifundefined{mathindent}%
  {\newdimen\mathindent\mathindent\leftmargini}%
  {}%

\def\resethooks{%
  \global\let\SaveRestoreHook\empty
  \global\let\ColumnHook\empty}
\newcommand*{\savecolumns}[1][default]%
  {\g@addto@macro\SaveRestoreHook{\savecolumns[#1]}}
\newcommand*{\restorecolumns}[1][default]%
  {\g@addto@macro\SaveRestoreHook{\restorecolumns[#1]}}
\newcommand*{\aligncolumn}[2]%
  {\g@addto@macro\ColumnHook{\column{#1}{#2}}}

\resethooks

\newcommand{\onelinecommentchars}{\quad-{}- }
\newcommand{\commentbeginchars}{\enskip\{-}
\newcommand{\commentendchars}{-\}\enskip}

\newcommand{\visiblecomments}{%
  \let\onelinecomment=\onelinecommentchars
  \let\commentbegin=\commentbeginchars
  \let\commentend=\commentendchars}

\newcommand{\invisiblecomments}{%
  \let\onelinecomment=\empty
  \let\commentbegin=\empty
  \let\commentend=\empty}

\visiblecomments

\newlength{\blanklineskip}
\setlength{\blanklineskip}{0.66084ex}

\newcommand{\hsindent}[1]{\quad}% default is fixed indentation
\let\hspre\empty
\let\hspost\empty
\newcommand{\NB}{\textbf{NB}}
\newcommand{\Todo}[1]{$\langle$\textbf{To do:}~#1$\rangle$}

\EndFmtInput
\makeatother
%
%
%
%
%
%
% This package provides two environments suitable to take the place
% of hscode, called "plainhscode" and "arrayhscode". 
%
% The plain environment surrounds each code block by vertical space,
% and it uses \abovedisplayskip and \belowdisplayskip to get spacing
% similar to formulas. Note that if these dimensions are changed,
% the spacing around displayed math formulas changes as well.
% All code is indented using \leftskip.
%
% Changed 19.08.2004 to reflect changes in colorcode. Should work with
% CodeGroup.sty.
%
\ReadOnlyOnce{polycode.fmt}%
\makeatletter

\newcommand{\hsnewpar}[1]%
  {{\parskip=0pt\parindent=0pt\par\vskip #1\noindent}}

% can be used, for instance, to redefine the code size, by setting the
% command to \small or something alike
\newcommand{\hscodestyle}{}

% The command \sethscode can be used to switch the code formatting
% behaviour by mapping the hscode environment in the subst directive
% to a new LaTeX environment.

\newcommand{\sethscode}[1]%
  {\expandafter\let\expandafter\hscode\csname #1\endcsname
   \expandafter\let\expandafter\endhscode\csname end#1\endcsname}

% "compatibility" mode restores the non-polycode.fmt layout.

\newenvironment{compathscode}%
  {\par\noindent
   \advance\leftskip\mathindent
   \hscodestyle
   \let\\=\@normalcr
   \let\hspre\(\let\hspost\)%
   \pboxed}%
  {\endpboxed\)%
   \par\noindent
   \ignorespacesafterend}

\newcommand{\compaths}{\sethscode{compathscode}}

% "plain" mode is the proposed default.
% It should now work with \centering.
% This required some changes. The old version
% is still available for reference as oldplainhscode.

\newenvironment{plainhscode}%
  {\hsnewpar\abovedisplayskip
   \advance\leftskip\mathindent
   \hscodestyle
   \let\hspre\(\let\hspost\)%
   \pboxed}%
  {\endpboxed%
   \hsnewpar\belowdisplayskip
   \ignorespacesafterend}

\newenvironment{oldplainhscode}%
  {\hsnewpar\abovedisplayskip
   \advance\leftskip\mathindent
   \hscodestyle
   \let\\=\@normalcr
   \(\pboxed}%
  {\endpboxed\)%
   \hsnewpar\belowdisplayskip
   \ignorespacesafterend}

% Here, we make plainhscode the default environment.

\newcommand{\plainhs}{\sethscode{plainhscode}}
\newcommand{\oldplainhs}{\sethscode{oldplainhscode}}
\plainhs

% The arrayhscode is like plain, but makes use of polytable's
% parray environment which disallows page breaks in code blocks.

\newenvironment{arrayhscode}%
  {\hsnewpar\abovedisplayskip
   \advance\leftskip\mathindent
   \hscodestyle
   \let\\=\@normalcr
   \(\parray}%
  {\endparray\)%
   \hsnewpar\belowdisplayskip
   \ignorespacesafterend}

\newcommand{\arrayhs}{\sethscode{arrayhscode}}

% The mathhscode environment also makes use of polytable's parray 
% environment. It is supposed to be used only inside math mode 
% (I used it to typeset the type rules in my thesis).

\newenvironment{mathhscode}%
  {\parray}{\endparray}

\newcommand{\mathhs}{\sethscode{mathhscode}}

% texths is similar to mathhs, but works in text mode.

\newenvironment{texthscode}%
  {\(\parray}{\endparray\)}

\newcommand{\texths}{\sethscode{texthscode}}

% The framed environment places code in a framed box.

\def\codeframewidth{\arrayrulewidth}
\RequirePackage{calc}

\newenvironment{framedhscode}%
  {\parskip=\abovedisplayskip\par\noindent
   \hscodestyle
   \arrayrulewidth=\codeframewidth
   \tabular{@{}|p{\linewidth-2\arraycolsep-2\arrayrulewidth-2pt}|@{}}%
   \hline\framedhslinecorrect\\{-1.5ex}%
   \let\endoflinesave=\\
   \let\\=\@normalcr
   \(\pboxed}%
  {\endpboxed\)%
   \framedhslinecorrect\endoflinesave{.5ex}\hline
   \endtabular
   \parskip=\belowdisplayskip\par\noindent
   \ignorespacesafterend}

\newcommand{\framedhslinecorrect}[2]%
  {#1[#2]}

\newcommand{\framedhs}{\sethscode{framedhscode}}

% The inlinehscode environment is an experimental environment
% that can be used to typeset displayed code inline.

\newenvironment{inlinehscode}%
  {\(\def\column##1##2{}%
   \let\>\undefined\let\<\undefined\let\\\undefined
   \newcommand\>[1][]{}\newcommand\<[1][]{}\newcommand\\[1][]{}%
   \def\fromto##1##2##3{##3}%
   \def\nextline{}}{\) }%

\newcommand{\inlinehs}{\sethscode{inlinehscode}}

% The joincode environment is a separate environment that
% can be used to surround and thereby connect multiple code
% blocks.

\newenvironment{joincode}%
  {\let\orighscode=\hscode
   \let\origendhscode=\endhscode
   \def\endhscode{\def\hscode{\endgroup\def\@currenvir{hscode}\\}\begingroup}
   %\let\SaveRestoreHook=\empty
   %\let\ColumnHook=\empty
   %\let\resethooks=\empty
   \orighscode\def\hscode{\endgroup\def\@currenvir{hscode}}}%
  {\origendhscode
   \global\let\hscode=\orighscode
   \global\let\endhscode=\origendhscode}%

\makeatother
\EndFmtInput
%


% format { = "\lbag"
% format } = "\rbag"


% TODO: fix spacing for \lbag and \rbag





















































%formap phiX = "\Phi"













\section{Parsing and reflective printing}
\label{sec:BiYacc}

\ignore{
\begin{hscode}\SaveRestoreHook
\column{B}{@{}>{\hspre}l<{\hspost}@{}}%
\column{E}{@{}>{\hspre}l<{\hspost}@{}}%
\>[B]{}\mbox{\enskip\{-\# LANGUAGE TemplateHaskell, TypeFamilies  \#-\}\enskip}{}\<[E]%
\\[\blanklineskip]%
\>[B]{}\mathbf{import}\;\Conid{\Conid{Generics}.BiGUL}{}\<[E]%
\\
\>[B]{}\mathbf{import}\;\Conid{\Conid{Generics}.\Conid{BiGUL}.TH}{}\<[E]%
\\
\>[B]{}\mathbf{import}\;\Conid{\Conid{Generics}.\Conid{BiGUL}.Lib}{}\<[E]%
\\
\>[B]{}\mathbf{import}\;\Conid{\Conid{Generics}.\Conid{BiGUL}.Interpreter}{}\<[E]%
\\[\blanklineskip]%
\>[B]{}\mathbf{import}\;\Conid{\Conid{GHC}.Generics}{}\<[E]%
\\
\>[B]{}\mathbf{import}\;\Conid{\Conid{Control}.Monad}{}\<[E]%
\\
\>[B]{}\mathbf{import}\;\Conid{\Conid{Text}.\Conid{ParserCombinators}.Parsec}{}\<[E]%
\\
\>[B]{}\mathbf{import}\;\Conid{\Conid{Text}.\Conid{ParserCombinators}.\Conid{Parsec}.Pos}{}\<[E]%
\ColumnHook
\end{hscode}\resethooks
}

When we mention the \emph{front-end} of a compiler, we usually think of a \emph{parser} that turns \emph{concrete syntax}, which is designed to be programmer-friendly and provides convenient syntactic sugar, into \emph{abstract syntax}, which is concise, structured, and easily manipulable by the compiler back-end.
There is another direction, though, in which a \emph{printer} turns abstract syntax back into concrete syntax.
This is useful, for example, for reporting the result of compiler optimizations done on abstract syntax to the programmer, who knows only concrete syntax.
In this case, though, we would want to print the optimized program in a form that is as close to the original program as possible, so the programmer can spot what are changed --- and not changed --- correctly and more easily.
This is where the notion of \emph{reflective printing} comes in: By taking both the original concrete program and the optimized abstract program as input, we can try to retain the look of the original program as much as possible.
Below we will use a simplified arithmetic expression language to explain how reflective printing can be implemented in BiGUL.

\subsection{Well-behavedness}

It is probably obvious that the idea of reflective printing comes from \ensuremath{\Varid{put}} transformations; parsing, then, is the \ensuremath{\Varid{get}} direction.
Before we proceed to implement parsing and reflective printing in BiGUL, a natural question to ask is:
Is well-behavedness meaningful in the context of parsing of reflective printing?
The answer is yes, especially for \ref{eq:PutGet}: An abstract syntax tree (AST) may be thought of as a concise and canonical representation of a concrete program, so it would be strange if a concrete program printed from an AST could not be parsed back to the same AST.
\ref{eq:GetPut}, on the other hand, is in fact not strong enough for our purpose, as it only says that, when an AST is unmodified, printing it reflectively to the original program does not change anything, whereas we would have liked to also say that ``small'' changes to the AST leads to only ``small'' changes to the concrete program.
It is at least a good start to have \ref{eq:GetPut}, though.
We thus conclude that BiGUL is indeed a suitable language for implementing reflective printers and corresponding parsers.

\subsection{Additive expressions}

Here we use a minimal example which is simple and yet can demonstrate what reflective printing is capable of.
Consider the following abstract syntax of arithmetic expressions consisting of integer constant, addition, and subtraction:
\begin{hscode}\SaveRestoreHook
\column{B}{@{}>{\hspre}l<{\hspost}@{}}%
\column{3}{@{}>{\hspre}l<{\hspost}@{}}%
\column{13}{@{}>{\hspre}c<{\hspost}@{}}%
\column{13E}{@{}l@{}}%
\column{16}{@{}>{\hspre}l<{\hspost}@{}}%
\column{21}{@{}>{\hspre}l<{\hspost}@{}}%
\column{E}{@{}>{\hspre}l<{\hspost}@{}}%
\>[B]{}\mathbf{data}\;\Conid{Arith}{}\<[13]%
\>[13]{}\mathrel{=}{}\<[13E]%
\>[16]{}\Conid{Num}\;\mathit{Int}{}\<[E]%
\\
\>[13]{}\mid {}\<[13E]%
\>[16]{}\Conid{Add}\;{}\<[21]%
\>[21]{}\Conid{Arith}\;\Conid{Arith}{}\<[E]%
\\
\>[13]{}\mid {}\<[13E]%
\>[16]{}\Conid{Sub}\;{}\<[21]%
\>[21]{}\Conid{Arith}\;\Conid{Arith}{}\<[E]%
\\
\>[B]{}\hsindent{3}{}\<[3]%
\>[3]{}\mathbf{deriving}\;\Conid{Show}{}\<[E]%
\ColumnHook
\end{hscode}\resethooks
This is a nice representation for the compiler, but we cannot expect the programmer to write something like ``\ensuremath{\Conid{Sub}\;(\Conid{Num}\;\mathrm{1})\;(\Conid{Add}\;(\Conid{Num}\;\mathrm{2})\;(\Conid{Num}\;\mathrm{3}))}'', and should provide a concrete syntax so that they can write ``$1 - (2+3)$''.
Such a concrete syntax is usually defined in terms of a BNF grammar:
\begin{hscode}\SaveRestoreHook
\column{B}{@{}>{\hspre}l<{\hspost}@{}}%
\column{9}{@{}>{\hspre}c<{\hspost}@{}}%
\column{9E}{@{}l@{}}%
\column{13}{@{}>{\hspre}l<{\hspost}@{}}%
\column{E}{@{}>{\hspre}l<{\hspost}@{}}%
\>[B]{}\Conid{Exp}{}\<[9]%
\>[9]{}\to {}\<[9E]%
\>[13]{}\Conid{Exp}\;\text{\tt '+'}\;\Conid{Factor}{}\<[E]%
\\
\>[9]{}\mid {}\<[9E]%
\>[13]{}\Conid{Exp}\;\text{\tt '-'}\;\Conid{Factor}{}\<[E]%
\\
\>[9]{}\mid {}\<[9E]%
\>[13]{}\Conid{Factor}{}\<[E]%
\\[\blanklineskip]%
\>[B]{}\Conid{Factor}{}\<[9]%
\>[9]{}\to {}\<[9E]%
\>[13]{}\mathit{Int}{}\<[E]%
\\
\>[9]{}\mid {}\<[9E]%
\>[13]{}\text{\tt '-'}\;\Conid{Factor}{}\<[E]%
\\
\>[9]{}\mid {}\<[9E]%
\>[13]{}\text{\tt '('}\;\Conid{Exp}\;\text{\tt ')'}{}\<[E]%
\ColumnHook
\end{hscode}\resethooks
The two-level structure of \ensuremath{\Conid{Exp}} and \ensuremath{\Conid{Factor}} ensures that plus and minus associate to the left by default; to change association, we should use parentheses.
And, to spice up the problem a little, we allow minus to be used also as a negative sign, as specified by the second production rule for \ensuremath{\Conid{Factor}}.
BiGUL deals with structured data only, so we should represent a string generated using this grammar as a concrete syntax tree of the following type:
\begin{hscode}\SaveRestoreHook
\column{B}{@{}>{\hspre}l<{\hspost}@{}}%
\column{11}{@{}>{\hspre}c<{\hspost}@{}}%
\column{11E}{@{}l@{}}%
\column{14}{@{}>{\hspre}l<{\hspost}@{}}%
\column{17}{@{}>{\hspre}l<{\hspost}@{}}%
\column{21}{@{}>{\hspre}l<{\hspost}@{}}%
\column{22}{@{}>{\hspre}l<{\hspost}@{}}%
\column{E}{@{}>{\hspre}l<{\hspost}@{}}%
\>[B]{}\mathbf{data}\;\Conid{Exp}{}\<[11]%
\>[11]{}\mathrel{=}{}\<[11E]%
\>[14]{}\Conid{Plus}\;{}\<[21]%
\>[21]{}\Conid{Exp}\;\Conid{Factor}{}\<[E]%
\\
\>[11]{}\mid {}\<[11E]%
\>[14]{}\Conid{Minus}\;{}\<[21]%
\>[21]{}\Conid{Exp}\;\Conid{Factor}{}\<[E]%
\\
\>[11]{}\mid {}\<[11E]%
\>[14]{}\Conid{EF}\;\Conid{Factor}{}\<[E]%
\\
\>[11]{}\mid {}\<[11E]%
\>[14]{}\Conid{ENull}{}\<[E]%
\\[\blanklineskip]%
\>[B]{}\mathbf{data}\;\Conid{Factor}{}\<[14]%
\>[14]{}\mathrel{=}{}\<[17]%
\>[17]{}\Conid{Lit}\;{}\<[22]%
\>[22]{}\mathit{Int}{}\<[E]%
\\
\>[14]{}\mid {}\<[17]%
\>[17]{}\Conid{Neg}\;{}\<[22]%
\>[22]{}\Conid{Factor}{}\<[E]%
\\
\>[14]{}\mid {}\<[17]%
\>[17]{}\Conid{Paren}\;\Conid{Exp}{}\<[E]%
\\
\>[14]{}\mid {}\<[17]%
\>[17]{}\Conid{FNull}{}\<[E]%
\ColumnHook
\end{hscode}\resethooks
%
\ignore{%
\begin{hscode}\SaveRestoreHook
\column{B}{@{}>{\hspre}l<{\hspost}@{}}%
\column{E}{@{}>{\hspre}l<{\hspost}@{}}%
\>[B]{}\Varid{deriveBiGULGeneric}\;\text{\tt ''}\;\Conid{Arith}{}\<[E]%
\\
\>[B]{}\Varid{deriveBiGULGeneric}\;\text{\tt ''}\;\Conid{Exp}{}\<[E]%
\\
\>[B]{}\Varid{deriveBiGULGeneric}\;\text{\tt ''}\;\Conid{Factor}{}\<[E]%
\ColumnHook
\end{hscode}\resethooks
}%
Apart from the \ensuremath{\Conid{Null}} constructors, which are inserted to represent incomplete trees that can occur during reflective printing, these two datatypes are in direct correspondence with the grammar, so it is easy to recover the string from a concrete syntax tree:
\begin{hscode}\SaveRestoreHook
\column{B}{@{}>{\hspre}l<{\hspost}@{}}%
\column{3}{@{}>{\hspre}l<{\hspost}@{}}%
\column{11}{@{}>{\hspre}l<{\hspost}@{}}%
\column{18}{@{}>{\hspre}l<{\hspost}@{}}%
\column{21}{@{}>{\hspre}l<{\hspost}@{}}%
\column{22}{@{}>{\hspre}c<{\hspost}@{}}%
\column{22E}{@{}l@{}}%
\column{25}{@{}>{\hspre}c<{\hspost}@{}}%
\column{25E}{@{}l@{}}%
\column{26}{@{}>{\hspre}c<{\hspost}@{}}%
\column{26E}{@{}l@{}}%
\column{28}{@{}>{\hspre}l<{\hspost}@{}}%
\column{29}{@{}>{\hspre}l<{\hspost}@{}}%
\column{E}{@{}>{\hspre}l<{\hspost}@{}}%
\>[B]{}\mathbf{instance}\;\Conid{Show}\;\Conid{Exp}\;\mathbf{where}{}\<[E]%
\\
\>[B]{}\hsindent{3}{}\<[3]%
\>[3]{}\Varid{show}\;({}\<[11]%
\>[11]{}\Conid{Plus}\;{}\<[18]%
\>[18]{}\Varid{e}\;{}\<[21]%
\>[21]{}\Varid{f}){}\<[25]%
\>[25]{}\mathrel{=}{}\<[25E]%
\>[28]{}\Varid{show}\;\Varid{e}\plus \text{\tt \char34 +\char34}\plus \Varid{show}\;\Varid{f}{}\<[E]%
\\
\>[B]{}\hsindent{3}{}\<[3]%
\>[3]{}\Varid{show}\;({}\<[11]%
\>[11]{}\Conid{Minus}\;{}\<[18]%
\>[18]{}\Varid{e}\;{}\<[21]%
\>[21]{}\Varid{f}){}\<[25]%
\>[25]{}\mathrel{=}{}\<[25E]%
\>[28]{}\Varid{show}\;\Varid{e}\plus \text{\tt \char34 -\char34}\plus \Varid{show}\;\Varid{f}{}\<[E]%
\\
\>[B]{}\hsindent{3}{}\<[3]%
\>[3]{}\Varid{show}\;({}\<[11]%
\>[11]{}\Conid{EF}\;{}\<[21]%
\>[21]{}\Varid{f}){}\<[25]%
\>[25]{}\mathrel{=}{}\<[25E]%
\>[28]{}\Varid{show}\;\Varid{f}{}\<[E]%
\\
\>[B]{}\hsindent{3}{}\<[3]%
\>[3]{}\Varid{show}\;{}\<[11]%
\>[11]{}\Conid{ENull}{}\<[25]%
\>[25]{}\mathrel{=}{}\<[25E]%
\>[28]{}\text{\tt \char34 .\char34}{}\<[E]%
\\[\blanklineskip]%
\>[B]{}\mathbf{instance}\;\Conid{Show}\;\Conid{Factor}\;\mathbf{where}{}\<[E]%
\\
\>[B]{}\hsindent{3}{}\<[3]%
\>[3]{}\Varid{show}\;({}\<[11]%
\>[11]{}\Conid{Lit}\;{}\<[18]%
\>[18]{}\Varid{n}{}\<[22]%
\>[22]{}){}\<[22E]%
\>[26]{}\mathrel{=}{}\<[26E]%
\>[29]{}\Varid{show}\;\Varid{n}{}\<[E]%
\\
\>[B]{}\hsindent{3}{}\<[3]%
\>[3]{}\Varid{show}\;({}\<[11]%
\>[11]{}\Conid{Neg}\;{}\<[18]%
\>[18]{}\Varid{f}{}\<[22]%
\>[22]{}){}\<[22E]%
\>[26]{}\mathrel{=}{}\<[26E]%
\>[29]{}\text{\tt \char34 -\char34}\plus \Varid{show}\;\Varid{f}{}\<[E]%
\\
\>[B]{}\hsindent{3}{}\<[3]%
\>[3]{}\Varid{show}\;({}\<[11]%
\>[11]{}\Conid{Paren}\;{}\<[18]%
\>[18]{}\Varid{e}{}\<[22]%
\>[22]{}){}\<[22E]%
\>[26]{}\mathrel{=}{}\<[26E]%
\>[29]{}\text{\tt \char34 (\char34}\plus \Varid{show}\;\Varid{e}\plus \text{\tt \char34 )\char34}{}\<[E]%
\\
\>[B]{}\hsindent{3}{}\<[3]%
\>[3]{}\Varid{show}\;{}\<[11]%
\>[11]{}\Conid{FNull}{}\<[26]%
\>[26]{}\mathrel{=}{}\<[26E]%
\>[29]{}\text{\tt \char34 .\char34}{}\<[E]%
\ColumnHook
\end{hscode}\resethooks
Conversely, using modern parser technologies like Haskell's \texttt{parsec} parser combinator library, we can easily implement a ``concrete parser'' that turns a string into a concrete syntax tree:
\begin{hscode}\SaveRestoreHook
\column{B}{@{}>{\hspre}l<{\hspost}@{}}%
\column{19}{@{}>{\hspre}c<{\hspost}@{}}%
\column{19E}{@{}l@{}}%
\column{23}{@{}>{\hspre}l<{\hspost}@{}}%
\column{E}{@{}>{\hspre}l<{\hspost}@{}}%
\>[B]{}\Varid{parseExp}\mathbin{::}\Conid{String}\to \Conid{Either}\;\Conid{ParseError}\;\Conid{Exp}{}\<[E]%
\\[\blanklineskip]%
\>[B]{}\Varid{unsafeParseExp}{}\<[19]%
\>[19]{}\mathbin{::}{}\<[19E]%
\>[23]{}\Conid{String}\to \Conid{Exp}{}\<[E]%
\\
\>[B]{}\Varid{unsafeParseExp}\;\Varid{s}{}\<[19]%
\>[19]{}\mathrel{=}{}\<[19E]%
\>[23]{}\mathbf{let}\;(\Conid{Right}\;\Varid{e})\mathrel{=}\Varid{parseExp}\;\Varid{s}\;\mathbf{in}\;\Varid{e}{}\<[E]%
\ColumnHook
\end{hscode}\resethooks
%
\ignore{%
\begin{hscode}\SaveRestoreHook
\column{B}{@{}>{\hspre}l<{\hspost}@{}}%
\column{3}{@{}>{\hspre}l<{\hspost}@{}}%
\column{5}{@{}>{\hspre}l<{\hspost}@{}}%
\column{6}{@{}>{\hspre}l<{\hspost}@{}}%
\column{8}{@{}>{\hspre}l<{\hspost}@{}}%
\column{15}{@{}>{\hspre}l<{\hspost}@{}}%
\column{18}{@{}>{\hspre}l<{\hspost}@{}}%
\column{19}{@{}>{\hspre}c<{\hspost}@{}}%
\column{19E}{@{}l@{}}%
\column{23}{@{}>{\hspre}l<{\hspost}@{}}%
\column{27}{@{}>{\hspre}l<{\hspost}@{}}%
\column{29}{@{}>{\hspre}l<{\hspost}@{}}%
\column{44}{@{}>{\hspre}l<{\hspost}@{}}%
\column{E}{@{}>{\hspre}l<{\hspost}@{}}%
\>[B]{}(\mathbin{>|>=})\mathbin{::}\Conid{Monad}\;\Varid{m}\Rightarrow \Varid{m}\;\Varid{a}\to (\Varid{a}\to \Varid{m}\;\Varid{b})\to \Varid{m}\;\Varid{a}{}\<[E]%
\\
\>[B]{}(\mathbin{>|>=})\;\Varid{mx}\;\Varid{f}\mathrel{=}\Varid{mx}\bind \lambda \Varid{x}\to \Varid{f}\;\Varid{x}\sequ \Varid{return}\;\Varid{x}{}\<[E]%
\\[\blanklineskip]%
\>[B]{}(\mathbin{>|>})\mathbin{::}\Conid{Monad}\;\Varid{m}\Rightarrow \Varid{m}\;\Varid{a}\to \Varid{m}\;\Varid{b}\to \Varid{m}\;\Varid{a}{}\<[E]%
\\
\>[B]{}(\mathbin{>|>})\;\Varid{mx}\;\Varid{my}\mathrel{=}\Varid{mx}\mathbin{>|>=}\Varid{const}\;\Varid{my}{}\<[E]%
\\[\blanklineskip]%
\>[B]{}\Varid{alternatives}\mathbin{::}[\mskip1.5mu \Conid{GenParser}\;\Varid{tok}\;\Varid{st}\;\Varid{a}\mskip1.5mu]\to \Conid{GenParser}\;\Varid{tok}\;\Varid{st}\;\Varid{a}{}\<[E]%
\\
\>[B]{}\Varid{alternatives}\mathrel{=}\Varid{foldr1}\;((\mathbin{<|>})\mathbin{\circ}\Varid{try}){}\<[E]%
\\[\blanklineskip]%
\>[B]{}\mathbf{data}\;\Conid{ExpToken}\mathrel{=}\Conid{LitTok}\;\mathit{Int}\mid \Conid{PlusTok}\mid \Conid{MinusTok}\mid \Conid{LParen}\mid \Conid{RParen}\;\mathbf{deriving}\;(\Conid{Eq},\Conid{Show}){}\<[E]%
\\[\blanklineskip]%
\>[B]{}\Varid{expTokeniser}\mathbin{::}\Conid{Parser}\;\Conid{ExpToken}{}\<[E]%
\\
\>[B]{}\Varid{expTokeniser}\mathrel{=}{}\<[E]%
\\
\>[B]{}\hsindent{3}{}\<[3]%
\>[3]{}\Varid{alternatives}\;{}\<[E]%
\\
\>[3]{}\hsindent{2}{}\<[5]%
\>[5]{}[\mskip1.5mu \Varid{liftM}\;(\Conid{LitTok}\mathbin{\circ}\Varid{read})\;(\Varid{many1}\;\Varid{digit}),{}\<[E]%
\\
\>[5]{}\hsindent{1}{}\<[6]%
\>[6]{}\Varid{char}\;\text{\tt '+'}\sequ \Varid{return}\;\Conid{PlusTok},{}\<[E]%
\\
\>[5]{}\hsindent{1}{}\<[6]%
\>[6]{}\Varid{char}\;\text{\tt '-'}\sequ \Varid{return}\;\Conid{MinusTok},{}\<[E]%
\\
\>[5]{}\hsindent{1}{}\<[6]%
\>[6]{}\Varid{char}\;\text{\tt '('}\sequ \Varid{return}\;\Conid{LParen},{}\<[E]%
\\
\>[5]{}\hsindent{1}{}\<[6]%
\>[6]{}\Varid{char}\;\text{\tt ')'}\sequ \Varid{return}\;\Conid{RParen}\mskip1.5mu]{}\<[E]%
\\[\blanklineskip]%
\>[B]{}\Varid{lit}\mathbin{::}\Conid{GenParser}\;\Conid{ExpToken}\;()\;\mathit{Int}{}\<[E]%
\\
\>[B]{}\Varid{lit}\mathrel{=}\Varid{token}\;\Varid{show}\;(\Varid{const}\;(\Varid{initialPos}\;\text{\tt \char34 \char34}))\;(\lambda \Varid{t}\to \mathbf{case}\;\Varid{t}\;\mathbf{of}\;\{\Conid{LitTok}\;\Varid{n}\to \Conid{Just}\;\Varid{n};\anonymous \to \Conid{Nothing}\}){}\<[E]%
\\[\blanklineskip]%
\>[B]{}\Varid{expToken}\mathbin{::}\Conid{ExpToken}\to \Conid{GenParser}\;\Conid{ExpToken}\;()\;\Conid{ExpToken}{}\<[E]%
\\
\>[B]{}\Varid{expToken}\;\Varid{tok}\mathrel{=}\Varid{token}\;\Varid{show}\;(\Varid{const}\;(\Varid{initialPos}\;\text{\tt \char34 \char34}))\;(\lambda \Varid{t}\to \mathbf{if}\;\Varid{t}\equiv \Varid{tok}\;\mathbf{then}\;\Conid{Just}\;\Varid{tok}\;\mathbf{else}\;\Conid{Nothing}){}\<[E]%
\\[\blanklineskip]%
\>[B]{}\Varid{expParser}\mathbin{::}\Conid{GenParser}\;\Conid{ExpToken}\;()\;\Conid{Exp}{}\<[E]%
\\
\>[B]{}\Varid{expParser}\mathrel{=}{}\<[E]%
\\
\>[B]{}\hsindent{3}{}\<[3]%
\>[3]{}\Varid{liftM}\;(\Varid{either}\;\Conid{EF}\;\Varid{id})\;{}\<[E]%
\\
\>[3]{}\hsindent{2}{}\<[5]%
\>[5]{}(\Varid{chainl1}\;{}\<[E]%
\\
\>[5]{}\hsindent{3}{}\<[8]%
\>[8]{}(\Varid{liftM}\;\Conid{Left}\;\Varid{factorParser})\;{}\<[E]%
\\
\>[5]{}\hsindent{3}{}\<[8]%
\>[8]{}(\Varid{liftM}\;(\lambda \Varid{op}\to \Varid{f}\;(\mathbf{if}\;\Varid{op}\equiv \Conid{PlusTok}\;\mathbf{then}\;\Conid{Plus}\;\mathbf{else}\;\Conid{Minus}))\;{}\<[E]%
\\
\>[8]{}\hsindent{7}{}\<[15]%
\>[15]{}(\Varid{alternatives}\;[\mskip1.5mu \Varid{expToken}\;\Conid{PlusTok},\Varid{expToken}\;\Conid{MinusTok}\mskip1.5mu]))){}\<[E]%
\\
\>[B]{}\hsindent{3}{}\<[3]%
\>[3]{}\mathbf{where}{}\<[E]%
\\
\>[3]{}\hsindent{2}{}\<[5]%
\>[5]{}\Varid{f}\mathbin{::}(\Conid{Exp}\to \Conid{Factor}\to \Conid{Exp})\to \Conid{Either}\;\Conid{Factor}\;\Conid{Exp}\to \Conid{Either}\;\Conid{Factor}\;\Conid{Exp}\to \Conid{Either}\;\Conid{Factor}\;\Conid{Exp}{}\<[E]%
\\
\>[3]{}\hsindent{2}{}\<[5]%
\>[5]{}\Varid{f}\;\Varid{con}\;(\Conid{Left}\;{}\<[18]%
\>[18]{}\Varid{lhsFactor})\;(\Conid{Left}\;\Varid{rhsFactor})\mathrel{=}\Conid{Right}\;(\Varid{con}\;(\Conid{EF}\;\Varid{lhsFactor})\;\Varid{rhsFactor}){}\<[E]%
\\
\>[3]{}\hsindent{2}{}\<[5]%
\>[5]{}\Varid{f}\;\Varid{con}\;(\Conid{Right}\;\Varid{lhsExp}{}\<[27]%
\>[27]{})\;(\Conid{Left}\;\Varid{rhsFactor})\mathrel{=}\Conid{Right}\;(\Varid{con}\;\Varid{lhsExp}\;\Varid{rhsFactor}){}\<[E]%
\\
\>[3]{}\hsindent{2}{}\<[5]%
\>[5]{}\Varid{f}\;\Varid{con}\;\anonymous \;{}\<[29]%
\>[29]{}(\Conid{Right}\;\anonymous {}\<[44]%
\>[44]{})\mathrel{=}\Varid{error}\;\text{\tt \char34 expParser:~the~impossible~happened\char34}{}\<[E]%
\\[\blanklineskip]%
\>[B]{}\Varid{factorParser}\mathbin{::}\Conid{GenParser}\;\Conid{ExpToken}\;()\;\Conid{Factor}{}\<[E]%
\\
\>[B]{}\Varid{factorParser}\mathrel{=}{}\<[E]%
\\
\>[B]{}\hsindent{3}{}\<[3]%
\>[3]{}\Varid{alternatives}\;{}\<[E]%
\\
\>[3]{}\hsindent{2}{}\<[5]%
\>[5]{}[\mskip1.5mu \Varid{liftM}\;\Conid{Lit}\;\Varid{lit},{}\<[E]%
\\
\>[5]{}\hsindent{1}{}\<[6]%
\>[6]{}\Varid{expToken}\;\Conid{MinusTok}\sequ \Varid{liftM}\;\Conid{Neg}\;\Varid{factorParser},{}\<[E]%
\\
\>[5]{}\hsindent{1}{}\<[6]%
\>[6]{}\Varid{expToken}\;\Conid{LParen}\sequ (\Varid{liftM}\;\Conid{Paren}\;\Varid{expParser}\mathbin{>|>}\Varid{expToken}\;\Conid{RParen})\mskip1.5mu]{}\<[E]%
\\[\blanklineskip]%
\>[B]{}\Varid{tokeniseAndParse}\mathbin{::}\Conid{Parser}\;\Varid{tok}\to \Conid{GenParser}\;\Varid{tok}\;()\;\Varid{a}\to \Conid{String}\to \Conid{Either}\;\Conid{ParseError}\;\Varid{a}{}\<[E]%
\\
\>[B]{}\Varid{tokeniseAndParse}\;\Varid{tokeniser}\;\Varid{parser}\mathrel{=}(\Varid{parse}\;\Varid{parser}\;\text{\tt \char34 \char34}\rbind )\mathbin{\circ}\Varid{parse}\;(\Varid{many}\;\Varid{tokeniser}\mathbin{>|>}\Varid{eof})\;\text{\tt \char34 \char34}{}\<[E]%
\\[\blanklineskip]%
\>[B]{}\Varid{parseExp}\mathbin{::}\Conid{String}\to \Conid{Either}\;\Conid{ParseError}\;\Conid{Exp}{}\<[E]%
\\
\>[B]{}\Varid{parseExp}\mathrel{=}\Varid{tokeniseAndParse}\;\Varid{expTokeniser}\;\Varid{expParser}{}\<[E]%
\\[\blanklineskip]%
\>[B]{}\Varid{unsafeParseExp}{}\<[19]%
\>[19]{}\mathbin{::}{}\<[19E]%
\>[23]{}\Conid{String}\to \Conid{Exp}{}\<[E]%
\\
\>[B]{}\Varid{unsafeParseExp}\;\Varid{s}{}\<[19]%
\>[19]{}\mathrel{=}{}\<[19E]%
\>[23]{}\mathbf{let}\;(\Conid{Right}\;\Varid{e})\mathrel{=}\Varid{parseExp}\;\Varid{s}\;\mathbf{in}\;\Varid{e}{}\<[E]%
\ColumnHook
\end{hscode}\resethooks
}%
The rest of the job is then write a BiGUL program between \ensuremath{\Conid{Exp}} and \ensuremath{\Conid{Arith}}.

\subsection{Reflective printing in BiGUL}

The program is basically a case analysis: For example, when the concrete side is a plus and the abstract side is an addition, they match, and we can go into their sub-trees recursively.
For the concrete side, the right sub-tree is of type \ensuremath{\Conid{Factor}} instead of \ensuremath{\Conid{Exp}}, so in fact we will write two (mutually recursive) programs:
\begin{hscode}\SaveRestoreHook
\column{B}{@{}>{\hspre}l<{\hspost}@{}}%
\column{15}{@{}>{\hspre}c<{\hspost}@{}}%
\column{15E}{@{}l@{}}%
\column{19}{@{}>{\hspre}l<{\hspost}@{}}%
\column{E}{@{}>{\hspre}l<{\hspost}@{}}%
\>[B]{}\Varid{pExpArith}{}\<[15]%
\>[15]{}\mathbin{::}{}\<[15E]%
\>[19]{}\Conid{BiGUL}\;\Conid{Exp}\;\Conid{Arith}{}\<[E]%
\\
\>[B]{}\Varid{pExpArith}{}\<[15]%
\>[15]{}\mathrel{=}{}\<[15E]%
\>[19]{}\Conid{Case}\;\bot {}\<[E]%
\\[\blanklineskip]%
\>[B]{}\Varid{pFactorArith}{}\<[15]%
\>[15]{}\mathbin{::}{}\<[15E]%
\>[19]{}\Conid{BiGUL}\;\Conid{Factor}\;\Conid{Arith}{}\<[E]%
\\
\>[B]{}\Varid{pFactorArith}{}\<[15]%
\>[15]{}\mathrel{=}{}\<[15E]%
\>[19]{}\Conid{Case}\;\bot {}\<[E]%
\ColumnHook
\end{hscode}\resethooks
The branch for plus and addition can then be written as:
\begin{hscode}\SaveRestoreHook
\column{B}{@{}>{\hspre}l<{\hspost}@{}}%
\column{11}{@{}>{\hspre}l<{\hspost}@{}}%
\column{E}{@{}>{\hspre}l<{\hspost}@{}}%
\>[B]{}{^\$}(\Varid{update}\;{}\<[11]%
\>[11]{}\mathrlap{^\mathbb{P}}\phantom{^\mathbb{D}}[\kern-1.5pt[\,\Conid{Plus}\;\Varid{l}\;\Varid{r}\,]\kern-1.5pt]\;\mathrlap{^\mathbb{P}}\phantom{^\mathbb{D}}[\kern-1.5pt[\,\Conid{Add}\;\Varid{l}\;\Varid{r}\,]\kern-1.5pt]\;{}\<[E]%
\\
\>[11]{}{^\mathbb{D}}[\kern-1.5pt[\,\Varid{l}\mathrel{=}\Varid{pExpArith};\Varid{r}\mathrel{=}\Varid{pFactorArith}\,]\kern-1.5pt]){}\<[E]%
\ColumnHook
\end{hscode}\resethooks
Following the same line of thought, we can fill in other branches to relate all abstract constructors with concrete production rules:
\begin{hscode}\SaveRestoreHook
\column{B}{@{}>{\hspre}l<{\hspost}@{}}%
\column{3}{@{}>{\hspre}l<{\hspost}@{}}%
\column{5}{@{}>{\hspre}l<{\hspost}@{}}%
\column{12}{@{}>{\hspre}c<{\hspost}@{}}%
\column{12E}{@{}l@{}}%
\column{15}{@{}>{\hspre}c<{\hspost}@{}}%
\column{15E}{@{}l@{}}%
\column{16}{@{}>{\hspre}l<{\hspost}@{}}%
\column{19}{@{}>{\hspre}l<{\hspost}@{}}%
\column{E}{@{}>{\hspre}l<{\hspost}@{}}%
\>[B]{}\Varid{pExpArith}{}\<[12]%
\>[12]{}\mathbin{::}{}\<[12E]%
\>[16]{}\Conid{BiGUL}\;\Conid{Exp}\;\Conid{Arith}{}\<[E]%
\\
\>[B]{}\Varid{pExpArith}{}\<[12]%
\>[12]{}\mathrel{=}{}\<[12E]%
\>[16]{}\Conid{Case}{}\<[E]%
\\
\>[B]{}\hsindent{3}{}\<[3]%
\>[3]{}[\mskip1.5mu {^\$}(\Varid{normalSV}\;\mathrlap{^\mathbb{P}}\phantom{^\mathbb{D}}[\kern-1.5pt[\,\Conid{Plus}\;\anonymous \;\anonymous \,]\kern-1.5pt]\;\mathrlap{^\mathbb{P}}\phantom{^\mathbb{D}}[\kern-1.5pt[\,\Conid{Add}\;\anonymous \;\anonymous \,]\kern-1.5pt]\;\mathrlap{^\mathbb{P}}\phantom{^\mathbb{D}}[\kern-1.5pt[\,\Conid{Plus}\;\anonymous \;\anonymous \,]\kern-1.5pt]){}\<[E]%
\\
\>[3]{}\hsindent{2}{}\<[5]%
\>[5]{}\Longrightarrow{^\$}(\Varid{update}\;{}\<[19]%
\>[19]{}\mathrlap{^\mathbb{P}}\phantom{^\mathbb{D}}[\kern-1.5pt[\,\Conid{Plus}\;\Varid{l}\;\Varid{r}\,]\kern-1.5pt]\;\mathrlap{^\mathbb{P}}\phantom{^\mathbb{D}}[\kern-1.5pt[\,\Conid{Add}\;\Varid{l}\;\Varid{r}\,]\kern-1.5pt]\;{}\<[E]%
\\
\>[19]{}{^\mathbb{D}}[\kern-1.5pt[\,\Varid{l}\mathrel{=}\Varid{pExpArith};\Varid{r}\mathrel{=}\Varid{pFactorArith}\,]\kern-1.5pt]){}\<[E]%
\\
\>[B]{}\hsindent{3}{}\<[3]%
\>[3]{},{^\$}(\Varid{normalSV}\;\mathrlap{^\mathbb{P}}\phantom{^\mathbb{D}}[\kern-1.5pt[\,\Conid{Minus}\;\anonymous \;\anonymous \,]\kern-1.5pt]\;\mathrlap{^\mathbb{P}}\phantom{^\mathbb{D}}[\kern-1.5pt[\,\Conid{Sub}\;\anonymous \;\anonymous \,]\kern-1.5pt]\;\mathrlap{^\mathbb{P}}\phantom{^\mathbb{D}}[\kern-1.5pt[\,\Conid{Minus}\;\anonymous \;\anonymous \,]\kern-1.5pt]){}\<[E]%
\\
\>[3]{}\hsindent{2}{}\<[5]%
\>[5]{}\Longrightarrow{^\$}(\Varid{update}\;{}\<[19]%
\>[19]{}\mathrlap{^\mathbb{P}}\phantom{^\mathbb{D}}[\kern-1.5pt[\,\Conid{Minus}\;\Varid{l}\;\Varid{r}\,]\kern-1.5pt]\;\mathrlap{^\mathbb{P}}\phantom{^\mathbb{D}}[\kern-1.5pt[\,\Conid{Sub}\;\Varid{l}\;\Varid{r}\,]\kern-1.5pt]\;{}\<[E]%
\\
\>[19]{}{^\mathbb{D}}[\kern-1.5pt[\,\Varid{l}\mathrel{=}\Varid{pExpArith};\Varid{r}\mathrel{=}\Varid{pFactorArith}\,]\kern-1.5pt]){}\<[E]%
\\
\>[B]{}\hsindent{3}{}\<[3]%
\>[3]{},{^\$}(\Varid{normalSV}\;\mathrlap{^\mathbb{P}}\phantom{^\mathbb{D}}[\kern-1.5pt[\,\Conid{EF}\;\anonymous \,]\kern-1.5pt]\;\mathrlap{^\mathbb{P}}\phantom{^\mathbb{D}}[\kern-1.5pt[\,\anonymous \,]\kern-1.5pt]\;\mathrlap{^\mathbb{P}}\phantom{^\mathbb{D}}[\kern-1.5pt[\,\Conid{EF}\;\anonymous \,]\kern-1.5pt]){}\<[E]%
\\
\>[3]{}\hsindent{2}{}\<[5]%
\>[5]{}\Longrightarrow{^\$}(\Varid{update}\;{}\<[19]%
\>[19]{}\mathrlap{^\mathbb{P}}\phantom{^\mathbb{D}}[\kern-1.5pt[\,\Conid{EF}\;\Varid{t}\,]\kern-1.5pt]\;\mathrlap{^\mathbb{P}}\phantom{^\mathbb{D}}[\kern-1.5pt[\,\Varid{t}\,]\kern-1.5pt]\;{}\<[E]%
\\
\>[19]{}{^\mathbb{D}}[\kern-1.5pt[\,\Varid{t}\mathrel{=}\Varid{pFactorArith}\,]\kern-1.5pt]){}\<[E]%
\\
\>[B]{}\hsindent{3}{}\<[3]%
\>[3]{}\mskip1.5mu]{}\<[E]%
\\[\blanklineskip]%
\>[B]{}\Varid{pFactorArith}{}\<[15]%
\>[15]{}\mathbin{::}{}\<[15E]%
\>[19]{}\Conid{BiGUL}\;\Conid{Factor}\;\Conid{Arith}{}\<[E]%
\\
\>[B]{}\Varid{pFactorArith}{}\<[15]%
\>[15]{}\mathrel{=}{}\<[15E]%
\>[19]{}\Conid{Case}{}\<[E]%
\\
\>[B]{}\hsindent{3}{}\<[3]%
\>[3]{}[\mskip1.5mu {^\$}(\Varid{normalSV}\;\mathrlap{^\mathbb{P}}\phantom{^\mathbb{D}}[\kern-1.5pt[\,\Conid{Lit}\;\anonymous \,]\kern-1.5pt]\;\mathrlap{^\mathbb{P}}\phantom{^\mathbb{D}}[\kern-1.5pt[\,\Conid{Num}\;\anonymous \,]\kern-1.5pt]\;\mathrlap{^\mathbb{P}}\phantom{^\mathbb{D}}[\kern-1.5pt[\,\Conid{Lit}\;\anonymous \,]\kern-1.5pt]){}\<[E]%
\\
\>[3]{}\hsindent{2}{}\<[5]%
\>[5]{}\Longrightarrow{^\$}(\Varid{update}\;\mathrlap{^\mathbb{P}}\phantom{^\mathbb{D}}[\kern-1.5pt[\,\Conid{Lit}\;\Varid{i}\,]\kern-1.5pt]\;\mathrlap{^\mathbb{P}}\phantom{^\mathbb{D}}[\kern-1.5pt[\,\Conid{Num}\;\Varid{i}\,]\kern-1.5pt]\;{^\mathbb{D}}[\kern-1.5pt[\,\Varid{i}\mathrel{=}\Conid{Replace}\,]\kern-1.5pt]){}\<[E]%
\\
\>[B]{}\hsindent{3}{}\<[3]%
\>[3]{},{^\$}(\Varid{normalSV}\;\mathrlap{^\mathbb{P}}\phantom{^\mathbb{D}}[\kern-1.5pt[\,\Conid{Neg}\;\anonymous \,]\kern-1.5pt]\;\mathrlap{^\mathbb{P}}\phantom{^\mathbb{D}}[\kern-1.5pt[\,\Conid{Sub}\;(\Conid{Num}\;\mathrm{0})\;\anonymous \,]\kern-1.5pt]\;\mathrlap{^\mathbb{P}}\phantom{^\mathbb{D}}[\kern-1.5pt[\,\Conid{Neg}\;\anonymous \,]\kern-1.5pt]){}\<[E]%
\\
\>[3]{}\hsindent{2}{}\<[5]%
\>[5]{}\Longrightarrow{^\$}(\Varid{update}\;\mathrlap{^\mathbb{P}}\phantom{^\mathbb{D}}[\kern-1.5pt[\,\Conid{Neg}\;\Varid{t}\,]\kern-1.5pt]\;\mathrlap{^\mathbb{P}}\phantom{^\mathbb{D}}[\kern-1.5pt[\,\Conid{Sub}\;(\Conid{Num}\;\mathrm{0})\;\Varid{t}\,]\kern-1.5pt]\;{^\mathbb{D}}[\kern-1.5pt[\,\Varid{t}\mathrel{=}\Varid{pFactorArith}\,]\kern-1.5pt]){}\<[E]%
\\
\>[B]{}\hsindent{3}{}\<[3]%
\>[3]{},{^\$}(\Varid{normalSV}\;\mathrlap{^\mathbb{P}}\phantom{^\mathbb{D}}[\kern-1.5pt[\,\Conid{Paren}\;\anonymous \,]\kern-1.5pt]\;\mathrlap{^\mathbb{P}}\phantom{^\mathbb{D}}[\kern-1.5pt[\,\anonymous \,]\kern-1.5pt]\;\mathrlap{^\mathbb{P}}\phantom{^\mathbb{D}}[\kern-1.5pt[\,\Conid{Paren}\;\anonymous \,]\kern-1.5pt]){}\<[E]%
\\
\>[3]{}\hsindent{2}{}\<[5]%
\>[5]{}\Longrightarrow{^\$}(\Varid{update}\;\mathrlap{^\mathbb{P}}\phantom{^\mathbb{D}}[\kern-1.5pt[\,\Conid{Paren}\;\Varid{t}\,]\kern-1.5pt]\;\mathrlap{^\mathbb{P}}\phantom{^\mathbb{D}}[\kern-1.5pt[\,\Varid{t}\,]\kern-1.5pt]\;{^\mathbb{D}}[\kern-1.5pt[\,\Varid{t}\mathrel{=}\Varid{pExpArith}\,]\kern-1.5pt]){}\<[E]%
\\
\>[B]{}\hsindent{3}{}\<[3]%
\>[3]{}\mskip1.5mu]{}\<[E]%
\ColumnHook
\end{hscode}\resethooks

This covers only ``normal'' cases though, namely when the source and view are ``the same'' except for parentheses and literals.
What about the cases when the source and view have mismatched shapes?
For these cases, we need adaptation.
Corresponding to each branch we have already written, we add an adaptive branch which looks at the shape of the view only, throws away a mismatched source, and creates an incomplete one whose shape matches that of the view; the source will be completely created through recursive processing.
For example, corresponding to the plus/addition branch, we write:
\begin{hscode}\SaveRestoreHook
\column{B}{@{}>{\hspre}l<{\hspost}@{}}%
\column{3}{@{}>{\hspre}l<{\hspost}@{}}%
\column{E}{@{}>{\hspre}l<{\hspost}@{}}%
\>[B]{}{^\$}(\Varid{adaptiveSV}\;\mathrlap{^\mathbb{P}}\phantom{^\mathbb{D}}[\kern-1.5pt[\,\anonymous \,]\kern-1.5pt]\;\mathrlap{^\mathbb{P}}\phantom{^\mathbb{D}}[\kern-1.5pt[\,\Conid{Add}\;\anonymous \;\anonymous \,]\kern-1.5pt]){}\<[E]%
\\
\>[B]{}\hsindent{3}{}\<[3]%
\>[3]{}\Longrightarrow\lambda \anonymous \;\anonymous \to \Conid{Plus}\;\Conid{ENull}\;\Conid{FNull}{}\<[E]%
\ColumnHook
\end{hscode}\resethooks
The full programs are:
\begin{hscode}\SaveRestoreHook
\column{B}{@{}>{\hspre}l<{\hspost}@{}}%
\column{3}{@{}>{\hspre}l<{\hspost}@{}}%
\column{5}{@{}>{\hspre}l<{\hspost}@{}}%
\column{12}{@{}>{\hspre}c<{\hspost}@{}}%
\column{12E}{@{}l@{}}%
\column{15}{@{}>{\hspre}c<{\hspost}@{}}%
\column{15E}{@{}l@{}}%
\column{16}{@{}>{\hspre}l<{\hspost}@{}}%
\column{18}{@{}>{\hspre}l<{\hspost}@{}}%
\column{19}{@{}>{\hspre}l<{\hspost}@{}}%
\column{E}{@{}>{\hspre}l<{\hspost}@{}}%
\>[B]{}\Varid{pExpArith}{}\<[12]%
\>[12]{}\mathbin{::}{}\<[12E]%
\>[16]{}\Conid{BiGUL}\;\Conid{Exp}\;\Conid{Arith}{}\<[E]%
\\
\>[B]{}\Varid{pExpArith}{}\<[12]%
\>[12]{}\mathrel{=}{}\<[12E]%
\>[16]{}\Conid{Case}{}\<[E]%
\\
\>[B]{}\hsindent{3}{}\<[3]%
\>[3]{}[\mskip1.5mu {^\$}(\Varid{normalSV}\;\mathrlap{^\mathbb{P}}\phantom{^\mathbb{D}}[\kern-1.5pt[\,\Conid{Plus}\;\anonymous \;\anonymous \,]\kern-1.5pt]\;\mathrlap{^\mathbb{P}}\phantom{^\mathbb{D}}[\kern-1.5pt[\,\Conid{Add}\;\anonymous \;\anonymous \,]\kern-1.5pt]\;\mathrlap{^\mathbb{P}}\phantom{^\mathbb{D}}[\kern-1.5pt[\,\Conid{Plus}\;\anonymous \;\anonymous \,]\kern-1.5pt]){}\<[E]%
\\
\>[3]{}\hsindent{2}{}\<[5]%
\>[5]{}\Longrightarrow{^\$}(\Varid{update}\;\mathrlap{^\mathbb{P}}\phantom{^\mathbb{D}}[\kern-1.5pt[\,\Conid{Plus}\;\Varid{l}\;\Varid{r}\,]\kern-1.5pt]\;\mathrlap{^\mathbb{P}}\phantom{^\mathbb{D}}[\kern-1.5pt[\,\Conid{Add}\;\Varid{l}\;\Varid{r}\,]\kern-1.5pt]\;{}\<[E]%
\\
\>[5]{}\hsindent{13}{}\<[18]%
\>[18]{}{^\mathbb{D}}[\kern-1.5pt[\,\Varid{l}\mathrel{=}\Varid{pExpArith};\Varid{r}\mathrel{=}\Varid{pFactorArith}\,]\kern-1.5pt]){}\<[E]%
\\
\>[B]{}\hsindent{3}{}\<[3]%
\>[3]{},{^\$}(\Varid{normalSV}\;\mathrlap{^\mathbb{P}}\phantom{^\mathbb{D}}[\kern-1.5pt[\,\Conid{Minus}\;\anonymous \;\anonymous \,]\kern-1.5pt]\;\mathrlap{^\mathbb{P}}\phantom{^\mathbb{D}}[\kern-1.5pt[\,\Conid{Sub}\;\anonymous \;\anonymous \,]\kern-1.5pt]\;\mathrlap{^\mathbb{P}}\phantom{^\mathbb{D}}[\kern-1.5pt[\,\Conid{Minus}\;\anonymous \;\anonymous \,]\kern-1.5pt]){}\<[E]%
\\
\>[3]{}\hsindent{2}{}\<[5]%
\>[5]{}\Longrightarrow{^\$}(\Varid{update}\;\mathrlap{^\mathbb{P}}\phantom{^\mathbb{D}}[\kern-1.5pt[\,\Conid{Minus}\;\Varid{l}\;\Varid{r}\,]\kern-1.5pt]\;\mathrlap{^\mathbb{P}}\phantom{^\mathbb{D}}[\kern-1.5pt[\,\Conid{Sub}\;\Varid{l}\;\Varid{r}\,]\kern-1.5pt]\;{}\<[E]%
\\
\>[5]{}\hsindent{13}{}\<[18]%
\>[18]{}{^\mathbb{D}}[\kern-1.5pt[\,\Varid{l}\mathrel{=}\Varid{pExpArith};\Varid{r}\mathrel{=}\Varid{pFactorArith}\,]\kern-1.5pt]){}\<[E]%
\\
\>[B]{}\hsindent{3}{}\<[3]%
\>[3]{},{^\$}(\Varid{normalSV}\;\mathrlap{^\mathbb{P}}\phantom{^\mathbb{D}}[\kern-1.5pt[\,\Conid{EF}\;\anonymous \,]\kern-1.5pt]\;\mathrlap{^\mathbb{P}}\phantom{^\mathbb{D}}[\kern-1.5pt[\,\anonymous \,]\kern-1.5pt]\;\mathrlap{^\mathbb{P}}\phantom{^\mathbb{D}}[\kern-1.5pt[\,\Conid{EF}\;\anonymous \,]\kern-1.5pt]){}\<[E]%
\\
\>[3]{}\hsindent{2}{}\<[5]%
\>[5]{}\Longrightarrow{^\$}(\Varid{update}\;\mathrlap{^\mathbb{P}}\phantom{^\mathbb{D}}[\kern-1.5pt[\,\Conid{EF}\;\Varid{t}\,]\kern-1.5pt]\;\mathrlap{^\mathbb{P}}\phantom{^\mathbb{D}}[\kern-1.5pt[\,\Varid{t}\,]\kern-1.5pt]\;{}\<[E]%
\\
\>[5]{}\hsindent{13}{}\<[18]%
\>[18]{}{^\mathbb{D}}[\kern-1.5pt[\,\Varid{t}\mathrel{=}\Varid{pFactorArith}\,]\kern-1.5pt]){}\<[E]%
\\
\>[B]{}\hsindent{3}{}\<[3]%
\>[3]{},{^\$}(\Varid{adaptiveSV}\;\mathrlap{^\mathbb{P}}\phantom{^\mathbb{D}}[\kern-1.5pt[\,\anonymous \,]\kern-1.5pt]\;\mathrlap{^\mathbb{P}}\phantom{^\mathbb{D}}[\kern-1.5pt[\,\Conid{Add}\;\anonymous \;\anonymous \,]\kern-1.5pt]){}\<[E]%
\\
\>[3]{}\hsindent{2}{}\<[5]%
\>[5]{}\Longrightarrow\lambda \anonymous \;\anonymous \to \Conid{Plus}\;\Conid{ENull}\;\Conid{FNull}{}\<[E]%
\\
\>[B]{}\hsindent{3}{}\<[3]%
\>[3]{},{^\$}(\Varid{adaptiveSV}\;\mathrlap{^\mathbb{P}}\phantom{^\mathbb{D}}[\kern-1.5pt[\,\anonymous \,]\kern-1.5pt]\;\mathrlap{^\mathbb{P}}\phantom{^\mathbb{D}}[\kern-1.5pt[\,\Conid{Sub}\;\anonymous \;\anonymous \,]\kern-1.5pt]){}\<[E]%
\\
\>[3]{}\hsindent{2}{}\<[5]%
\>[5]{}\Longrightarrow\lambda \anonymous \;\anonymous \to \Conid{Minus}\;\Conid{ENull}\;\Conid{FNull}{}\<[E]%
\\
\>[B]{}\hsindent{3}{}\<[3]%
\>[3]{},{^\$}(\Varid{adaptiveSV}\;\mathrlap{^\mathbb{P}}\phantom{^\mathbb{D}}[\kern-1.5pt[\,\anonymous \,]\kern-1.5pt]\;\mathrlap{^\mathbb{P}}\phantom{^\mathbb{D}}[\kern-1.5pt[\,\anonymous \,]\kern-1.5pt]){}\<[E]%
\\
\>[3]{}\hsindent{2}{}\<[5]%
\>[5]{}\Longrightarrow\lambda \anonymous \;\anonymous \to \Conid{EF}\;\Conid{FNull}{}\<[E]%
\\
\>[B]{}\hsindent{3}{}\<[3]%
\>[3]{}\mskip1.5mu]{}\<[E]%
\\[\blanklineskip]%
\>[B]{}\Varid{pFactorArith}{}\<[15]%
\>[15]{}\mathbin{::}{}\<[15E]%
\>[19]{}\Conid{BiGUL}\;\Conid{Factor}\;\Conid{Arith}{}\<[E]%
\\
\>[B]{}\Varid{pFactorArith}{}\<[15]%
\>[15]{}\mathrel{=}{}\<[15E]%
\>[19]{}\Conid{Case}{}\<[E]%
\\
\>[B]{}\hsindent{3}{}\<[3]%
\>[3]{}[\mskip1.5mu {^\$}(\Varid{normalSV}\;\mathrlap{^\mathbb{P}}\phantom{^\mathbb{D}}[\kern-1.5pt[\,\Conid{Lit}\;\anonymous \,]\kern-1.5pt]\;\mathrlap{^\mathbb{P}}\phantom{^\mathbb{D}}[\kern-1.5pt[\,\Conid{Num}\;\anonymous \,]\kern-1.5pt]\;\mathrlap{^\mathbb{P}}\phantom{^\mathbb{D}}[\kern-1.5pt[\,\Conid{Lit}\;\anonymous \,]\kern-1.5pt]){}\<[E]%
\\
\>[3]{}\hsindent{2}{}\<[5]%
\>[5]{}\Longrightarrow{^\$}(\Varid{update}\;\mathrlap{^\mathbb{P}}\phantom{^\mathbb{D}}[\kern-1.5pt[\,\Conid{Lit}\;\Varid{i}\,]\kern-1.5pt]\;\mathrlap{^\mathbb{P}}\phantom{^\mathbb{D}}[\kern-1.5pt[\,\Conid{Num}\;\Varid{i}\,]\kern-1.5pt]\;{^\mathbb{D}}[\kern-1.5pt[\,\Varid{i}\mathrel{=}\Conid{Replace}\,]\kern-1.5pt]){}\<[E]%
\\
\>[B]{}\hsindent{3}{}\<[3]%
\>[3]{},{^\$}(\Varid{normalSV}\;\mathrlap{^\mathbb{P}}\phantom{^\mathbb{D}}[\kern-1.5pt[\,\Conid{Neg}\;\anonymous \,]\kern-1.5pt]\;\mathrlap{^\mathbb{P}}\phantom{^\mathbb{D}}[\kern-1.5pt[\,\Conid{Sub}\;(\Conid{Num}\;\mathrm{0})\;\anonymous \,]\kern-1.5pt]\;\mathrlap{^\mathbb{P}}\phantom{^\mathbb{D}}[\kern-1.5pt[\,\Conid{Neg}\;\anonymous \,]\kern-1.5pt]){}\<[E]%
\\
\>[3]{}\hsindent{2}{}\<[5]%
\>[5]{}\Longrightarrow{^\$}(\Varid{update}\;\mathrlap{^\mathbb{P}}\phantom{^\mathbb{D}}[\kern-1.5pt[\,\Conid{Neg}\;\Varid{t}\,]\kern-1.5pt]\;\mathrlap{^\mathbb{P}}\phantom{^\mathbb{D}}[\kern-1.5pt[\,\Conid{Sub}\;(\Conid{Num}\;\mathrm{0})\;\Varid{t}\,]\kern-1.5pt]\;{^\mathbb{D}}[\kern-1.5pt[\,\Varid{t}\mathrel{=}\Varid{pFactorArith}\,]\kern-1.5pt]){}\<[E]%
\\
\>[B]{}\hsindent{3}{}\<[3]%
\>[3]{},{^\$}(\Varid{normalSV}\;\mathrlap{^\mathbb{P}}\phantom{^\mathbb{D}}[\kern-1.5pt[\,\Conid{Paren}\;\anonymous \,]\kern-1.5pt]\;\mathrlap{^\mathbb{P}}\phantom{^\mathbb{D}}[\kern-1.5pt[\,\anonymous \,]\kern-1.5pt]\;\mathrlap{^\mathbb{P}}\phantom{^\mathbb{D}}[\kern-1.5pt[\,\Conid{Paren}\;\anonymous \,]\kern-1.5pt]){}\<[E]%
\\
\>[3]{}\hsindent{2}{}\<[5]%
\>[5]{}\Longrightarrow{^\$}(\Varid{update}\;\mathrlap{^\mathbb{P}}\phantom{^\mathbb{D}}[\kern-1.5pt[\,\Conid{Paren}\;\Varid{t}\,]\kern-1.5pt]\;\mathrlap{^\mathbb{P}}\phantom{^\mathbb{D}}[\kern-1.5pt[\,\Varid{t}\,]\kern-1.5pt]\;{^\mathbb{D}}[\kern-1.5pt[\,\Varid{t}\mathrel{=}\Varid{pExpArith}\,]\kern-1.5pt]){}\<[E]%
\\
\>[B]{}\hsindent{3}{}\<[3]%
\>[3]{},{^\$}(\Varid{adaptiveSV}\;\mathrlap{^\mathbb{P}}\phantom{^\mathbb{D}}[\kern-1.5pt[\,\anonymous \,]\kern-1.5pt]\;\mathrlap{^\mathbb{P}}\phantom{^\mathbb{D}}[\kern-1.5pt[\,\Conid{Num}\;\anonymous \,]\kern-1.5pt]){}\<[E]%
\\
\>[3]{}\hsindent{2}{}\<[5]%
\>[5]{}\Longrightarrow\lambda \anonymous \;\anonymous \to \Conid{Lit}\;\mathrm{0}{}\<[E]%
\\
\>[B]{}\hsindent{3}{}\<[3]%
\>[3]{},{^\$}(\Varid{adaptiveSV}\;\mathrlap{^\mathbb{P}}\phantom{^\mathbb{D}}[\kern-1.5pt[\,\anonymous \,]\kern-1.5pt]\;\mathrlap{^\mathbb{P}}\phantom{^\mathbb{D}}[\kern-1.5pt[\,\Conid{Sub}\;(\Conid{Num}\;\mathrm{0})\;\anonymous \,]\kern-1.5pt]){}\<[E]%
\\
\>[3]{}\hsindent{2}{}\<[5]%
\>[5]{}\Longrightarrow\lambda \anonymous \;\anonymous \to \Conid{Neg}\;\Conid{FNull}{}\<[E]%
\\
\>[B]{}\hsindent{3}{}\<[3]%
\>[3]{},{^\$}(\Varid{adaptiveSV}\;\mathrlap{^\mathbb{P}}\phantom{^\mathbb{D}}[\kern-1.5pt[\,\anonymous \,]\kern-1.5pt]\;\mathrlap{^\mathbb{P}}\phantom{^\mathbb{D}}[\kern-1.5pt[\,\anonymous \,]\kern-1.5pt]){}\<[E]%
\\
\>[3]{}\hsindent{2}{}\<[5]%
\>[5]{}\Longrightarrow\lambda \anonymous \;\anonymous \to \Conid{Paren}\;\Conid{ENull}{}\<[E]%
\\
\>[B]{}\hsindent{3}{}\<[3]%
\>[3]{}\mskip1.5mu]{}\<[E]%
\ColumnHook
\end{hscode}\resethooks

\subsection{Reflecting optimizations and evaluation sequences}

The BiGUL programs, being bidirectional, can be executed in the \ensuremath{\Varid{put}} direction as a reflective printer, or in the \ensuremath{\Varid{get}} direction as a parser.
Let us look at parsing first. For example:
\begin{tabbing}\tt
~\char42{}Main\char62{}~get~pExpArith~\char40{}unsafeParseExp~\char34{}\char40{}\char45{}\char40{}3\char43{}4\char41{}\char41{}\char34{}\char41{}\\
\tt ~Just~\char40{}Sub~\char40{}Num~0\char41{}~\char40{}Add~\char40{}Num~3\char41{}~\char40{}Num~4\char41{}\char41{}\char41{}
\end{tabbing}
Note that a unary minus is considered as syntactic sugar, and is desugared into a subtraction whose left operand is zero.
Also note that parentheses are turned into correct structure of the abstract syntax tree, and nothing more --- excessive parentheses are cleanly discarded.

For reflective printing, as we mentioned, one application is reporting what compiler optimizations do.
We can replace the sub-expression $3+4$ with its value~$1$, for example, and the reflective printer will be able to retain the excessive parentheses:
\begin{tabbing}\tt
~\char42{}Main\char62{}~put~pExpArith~\char40{}unsafeParseExp~\char34{}\char40{}\char45{}\char40{}3\char43{}4\char41{}\char41{}\char34{}\char41{}\\
\tt ~~~~~~~~~~~~~~~~~~~~~~\char40{}Sub~\char40{}Num~0\char41{}~\char40{}Num~7\char41{}\char41{}\\
\tt ~Just~\char40{}\char45{}\char40{}7\char41{}\char41{}
\end{tabbing}
Notice also that the unary minus is preserved.
If the original concrete expression uses a binary minus instead, it will be preserved as well:
\begin{tabbing}\tt
~\char42{}Main\char62{}~put~pExpArith~\char40{}unsafeParseExp~\char34{}\char40{}0\char45{}\char40{}3\char43{}4\char41{}\char41{}\char34{}\char41{}\\
\tt ~~~~~~~~~~~~~~~~~~~~~~\char40{}Sub~\char40{}Num~0\char41{}~\char40{}Num~7\char41{}\char41{}\\
\tt ~Just~\char40{}0\char45{}\char40{}7\char41{}\char41{}
\end{tabbing}

More generally, the steps in an evaluation sequence of an abstract syntax tree can all be reflected to concrete syntax.
For example, starting from:
\begin{tabbing}\tt
~\char42{}Main\char62{}~get~pExpArith~\char40{}unsafeParseExp~\char34{}1\char43{}\char40{}2\char43{}3\char41{}\char34{}\char41{}\\
\tt ~Just~\char40{}Add~\char40{}Num~1\char41{}~\char40{}Add~\char40{}Num~2\char41{}~\char40{}Num~3\char41{}\char41{}\char41{}
\end{tabbing}
it takes two steps to evaluate this expression:
\begin{tabbing}\tt
~\char42{}Main\char62{}~put~pExpArith~\char40{}unsafeParseExp~\char34{}1\char43{}\char40{}2\char43{}3\char41{}\char34{}\char41{}\\
\tt ~~~~~~~~~~~~~~~~~~~~~~\char40{}Add~\char40{}Num~1\char41{}~\char40{}Num~5\char41{}\char41{}\\
\tt ~Just~1\char43{}\char40{}5\char41{}\\
\tt ~\char42{}Main\char62{}~put~pExpArith~\char40{}unsafeParseExp~\char34{}1\char43{}\char40{}5\char41{}\char34{}\char41{}~\char40{}Num~6\char41{}\\
\tt ~Just~6
\end{tabbing}
This means that if we have an evaluator on the abstract syntax, we will automatically get an evaluator on the concrete syntax!

A reflective printer can also be used as an ordinary printer by setting the original source to an empty one.
For example:
\begin{tabbing}\tt
~\char42{}Main\char62{}~put~pExpArith~ENull~\char40{}Sub~\char40{}Num~0\char41{}~\char40{}Add~\char40{}Num~1\char41{}~\char40{}Num~1\char41{}\char41{}\char41{}\\
\tt ~Just~0\char45{}\char40{}1\char43{}1\char41{}
\end{tabbing}
You have probably noticed that the subtraction is reflected as a binary minus instead of a unary one, despite that the left operand is zero.
This behavior is easily customizable:
By adding an adaptive branch before the one dealing generically with \ensuremath{\Conid{Sub}} in \ensuremath{\Varid{pExpArith}}:
\begin{hscode}\SaveRestoreHook
\column{B}{@{}>{\hspre}l<{\hspost}@{}}%
\column{3}{@{}>{\hspre}l<{\hspost}@{}}%
\column{E}{@{}>{\hspre}l<{\hspost}@{}}%
\>[B]{}{^\$}(\Varid{adaptiveSV}\;\mathrlap{^\mathbb{P}}\phantom{^\mathbb{D}}[\kern-1.5pt[\,\anonymous \,]\kern-1.5pt]\;\mathrlap{^\mathbb{P}}\phantom{^\mathbb{D}}[\kern-1.5pt[\,\Conid{Sub}\;(\Conid{Num}\;\mathrm{0})\;\anonymous \,]\kern-1.5pt]){}\<[E]%
\\
\>[B]{}\hsindent{3}{}\<[3]%
\>[3]{}\Longrightarrow\lambda \anonymous \;\anonymous \to \Conid{EF}\;\Conid{FNull}{}\<[E]%
\ColumnHook
\end{hscode}\resethooks
the above abstract syntax tree can be printed as:
\begin{tabbing}\tt
~\char42{}Main\char62{}~put~pExpArith~ENull~\char40{}Sub~\char40{}Num~0\char41{}~\char40{}Add~\char40{}Num~1\char41{}~\char40{}Num~1\char41{}\char41{}\char41{}\\
\tt ~Just~\char45{}\char40{}1\char43{}1\char41{}
\end{tabbing}

\subsection{A domain-specific language}

As a final remark, the above programs may look long, but at the core of them are merely the correspondences between concrete production rules and abstract constructors.
We can design a domain-specific language (DSL) that expresses such correspondences concisely, and then expand programs in this DSL into BiGUL.
In fact, we have already done so, and the DSL is called \emph{BiYacc}.
For example, all the programs we have written can be generated from the following eight-line BiYacc program:
\begin{tabbing}\tt
~~~Arith~\char43{}\char62{}~Exp\\
\tt ~~~Add~l~r~\char43{}\char62{}~\char40{}l~\char43{}\char62{}~Exp\char41{}~\char39{}\char43{}\char39{}~\char40{}r~\char43{}\char62{}~Factor\char41{}\char59{}\\
\tt ~~~Sub~l~r~\char43{}\char62{}~\char40{}l~\char43{}\char62{}~Exp\char41{}~\char39{}\char45{}\char39{}~\char40{}r~\char43{}\char62{}~Factor\char41{}\char59{}\\
\tt ~~~f~~~~~~~\char43{}\char62{}~\char40{}f~\char43{}\char62{}~Factor\char41{}\char59{}\\
\tt ~\\
\tt ~~~Arith~\char43{}\char62{}~Factor\\
\tt ~~~Num~n~~~~~~~~~\char43{}\char62{}~\char40{}n~\char43{}\char62{}~Int\char41{}\char59{}\\
\tt ~~~Sub~\char40{}Num~0\char41{}~r~\char43{}\char62{}~\char39{}\char45{}\char39{}~\char40{}r~\char43{}\char62{}~Factor\char41{}\char59{}\\
\tt ~~~f~~~~~~~~~~~~~\char43{}\char62{}~\char39{}\char40{}\char39{}~\char40{}f~\char43{}\char62{}~Exp\char41{}~\char39{}\char41{}\char39{}\char59{}
\end{tabbing}
See our SLE 2016 paper~\cite{Zhu-BiYacc} for more interesting experiments about reflective printing, done on a more realistic imperative language.
